\documentclass{beamer}
\usepackage{hyperref}
\usepackage{xcolor}
\usepackage{graphicx}
\usepackage{tikz}
\usetikzlibrary{positioning}
\usepackage{authoraftertitle}
\usepackage{fontspec}
\usepackage{listings}

% Setarea fonturilor specifice sablonului
\setsansfont{Oswald}[Path=./OswaldFontFiles/,Scale=0.9,Extension=.ttf,UprightFont=*-Regular,BoldFont=*-Bold]
\setromanfont{Lato}[Path=./LatoFontFiles/,Scale=0.9,Extension=.ttf,UprightFont=*-Regular,BoldFont=*-Bold]    

\setbeamertemplate{navigation symbols}{}
% Asigura-te ca imaginea de fundal exista in folderul img/
\setbeamertemplate{background}{\includegraphics[width=\paperwidth,height=\paperheight]{img/slider-in.jpg}}

% --- INFORMATIILE ECHIPE FATA ---
\author{S. Ioan-Victor, S. Hojda, T. Copaciu}
\title{Security Analysis of VulnBank Application}
% Folosim campurile sablonului pentru informatii despre curs
\newcommand\Facultate{Quality Aspects of Security in Software Testing}
\newcommand\Universitate{QASST Project 2026}
\newcommand\Adresa{Babe\c{s}-Bolyai University}
\newcommand\Website{CS Master - CyberSecurity}
\newcommand\DataPrezentarii{12/01/2026}

\usepackage{CSTheme_EN}

\definecolor{codegray}{rgb}{0.95,0.95,0.95}
\definecolor{codegreen}{rgb}{0,0.6,0}
\lstdefinestyle{mystyle}{
    backgroundcolor=\color{codegray},
    commentstyle=\color{codegreen},
    keywordstyle=\color{blue},
    basicstyle=\ttfamily\footnotesize,
    breakatwhitespace=false,
    breaklines=true,
    captionpos=b,
    keepspaces=true,
    showspaces=false,
    showstringspaces=false,
    showtabs=false,
    tabsize=2
}

\lstset{style=mystyle}

\begin{document}

% --- SLIDE 1: TITLE PAGE ---
% Title slide frame
\begin{frame}[plain]
\begin{tikzpicture}[remember picture,overlay]
% Background image
\node[above right,inner sep=0pt] at (current page.south west)
{
    \includegraphics[width=\paperwidth]{img/blueCS_background.png}
};
    
% Title & Subtitle

\node[above=1.5cm,color=yellowCS,align=center] (title) at (current page.center)
{
    \Huge \MyTitle
};

% Author 
\node[color=white,below=0.1cm] (author) at (title.south){\MyAuthor};
% Facultate
\node[color=yellowCS,below=1cm] (date) at (author.south){\large \Facultate};
% Universitate
\node[color=white,below=1.5cm] (date) at (author.south){\large \Universitate};
% Logo CS
\node[below right=3.5cm and 0.5cm] at (current page.north west)
{
    \includegraphics[width=2cm]{img/CS_white_logo.png}
};
\draw[yellowCS, ultra thick] (3,-0.65) -- (8,-0.65);
% Logo UBB
\node[below right =3.5cm and 10cm] at (current page.north west)
{
    \includegraphics[width=2cm]{img/Sigla_UBB_claudiopolitana_alb.png}
};
\end{tikzpicture}
    
\end{frame}


% --- SLIDE 2: AGENDA (Adaptat pe structura sablonului) ---
\newcommand\sectionTitle{Overview}
\newcommand\sectionSubtitle{Presentation Agenda}
\newcommand\sectionMotto{Software Tested, Strategy, Exploitation \& Remediation}
% Section 1 slide frame
\begin{frame}[plain]
\begin{tikzpicture}[remember picture,overlay]
% Background image
    \node[above right,inner sep=0pt] at (current page.south west)
    {
        \includegraphics[width=\paperwidth]{img/blueCS_background.png}
    };
% Logo CS 50
\node[below right=0.2cm and 7.5cm] at (current page.north west)
{
    \includegraphics[width=1.85cm]{img/facultatea de matematica_alba_transp.png}
};
% Logo CS
\node[below right=0.5cm and 10cm] at (current page.north west)
{
    \includegraphics[width=2cm]{img/CS_white_logo.png}
};
\draw[yellowCS, ultra thick] (0.1,-4) -- (6,-4);
  
% section title
\node[color=white,below right=2.5cm and 0cm] {\huge \sectionTitle};
\node[color=white,below right=3.2cm and 0cm] {\Large \sectionSubtitle};
\node[color=white,below left=4.7cm and 1.45cm] (author) at (title.south){ \sectionMotto};

\end{tikzpicture}
    
\end{frame}


\begin{frame}{Presentation Contents}
    \begin{itemize}
        \item \textbf{Target:} VulnBank Application Overview
        \item \textbf{Methodology:} Hybrid Approach (Defensive \& Offensive)
        \item \textbf{Findings:} Summary of Identified Vulnerabilities
        \item \textbf{Deep Dive:} Offensive Exploitation (SQL Injection)
        \item \textbf{Resolution:} Remediation Steps \& Conclusions
    \end{itemize}
\end{frame}

% --- SECTIUNEA 1: SOFTWARE TESTED ---
\renewcommand*\sectionTitle{Software Tested}
\renewcommand*\sectionSubtitle{VulnBank Application}
% Section 2 slide frame
\begin{frame}[plain]
\begin{tikzpicture}[remember picture,overlay]
% Background image
    \node[above right,inner sep=0pt] at (current page.south west)
    {
        \includegraphics[width=\paperwidth]{img/blueCS_background.png}
    };
% Logo CS
\node[below right =6cm and 0.5cm] at (current page.north west)
{
    \includegraphics[width=2cm]{img/CS_white_logo.png}
};
\draw[yellowCS, ultra thick] (2,-3.25) -- (8,-3.25);
  
% section title
\node[color=yellowCS,below right=2cm and 2cm] {\Huge \sectionTitle};
\node[color=white,below right=3.5cm and 2cm] {\Large \sectionSubtitle};

\end{tikzpicture}
    
\end{frame}


\begin{frame}{Target Application: VulnBank}
\begin{minipage}{0.55\linewidth}
        \begin{itemize}
            \item \textbf{Type:} Intentionally vulnerable banking simulation (Web \& API).
            \item \textbf{Critical Features:} 
            \begin{itemize}
                \item Fund Transfers
                \item Loan Requests
                \item Profile Management
            \end{itemize}
            \item \textbf{Why this target?} It handles sensitive financial data (PII, Transactions), making it an ideal candidate for security testing.
        \end{itemize}
 \end{minipage}
 \begin{minipage}{0.4\linewidth}
    \begin{figure}[h]
        \centering
        % Inlocuieste cu calea corecta a imaginii tale
        \includegraphics[width=\linewidth]{Figures/vuln_dashboard.png}
        \caption{VulnBank Dashboard}
    \end{figure}
\end{minipage}
\end{frame}

% --- SECTIUNEA 2: STRATEGY ---
\renewcommand*\sectionTitle{Security Approach}
\renewcommand*\sectionSubtitle{Hybrid Strategy}
% Section 2 slide frame
\begin{frame}[plain]
\begin{tikzpicture}[remember picture,overlay]
% Background image
    \node[above right,inner sep=0pt] at (current page.south west)
    {
        \includegraphics[width=\paperwidth]{img/blueCS_background.png}
    };
% Logo CS
\node[below right =6cm and 0.5cm] at (current page.north west)
{
    \includegraphics[width=2cm]{img/CS_white_logo.png}
};
\draw[yellowCS, ultra thick] (2,-3.25) -- (8,-3.25);
  
% section title
\node[color=yellowCS,below right=2cm and 2cm] {\Huge \sectionTitle};
\node[color=white,below right=3.5cm and 2cm] {\Large \sectionSubtitle};

\end{tikzpicture}
    
\end{frame}


\begin{frame}{Hybrid Strategy: White-box \& Black-box}
\begin{minipage}[t]{0.48\linewidth}
\begin{block}{\centering Defensive (White-box)}
        \begin{itemize}
            \item \textbf{Tool:} Snyk (SAST)
            \item \textbf{Focus:} Code Quality \& Configuration
            \item \textbf{Key Findings:}
            \begin{itemize}
                \item Insecure Code Patterns (f-strings)
                \item SSL Verification Disabled
                \item Missing Cookie Attributes
            \end{itemize}
        \end{itemize}
\end{block}        
 \end{minipage}
  \hspace{0.08\textwidth}%
 \begin{minipage}[t]{0.42\linewidth}
\begin{block}{\centering Offensive (Black-box)}        
        \begin{itemize}
            \item \textbf{Tools:} Burp Suite, Nessus (DAST)
            \item \textbf{Focus:} Runtime Exploitation
            \item \textbf{Key Findings:}
            \item \textbf{SQL Injection (API)}
            \item API Documentation Exposure
            \item Cleartext Credentials
        \end{itemize}
 \end{block}        
\end{minipage}
\end{frame}

% --- SECTIUNEA 3: FINDINGS & EXPLOIT ---
\renewcommand*\sectionTitle{Critical Findings}
\renewcommand*\sectionSubtitle{Offensive Exploitation}
% Section 2 slide frame
\begin{frame}[plain]
\begin{tikzpicture}[remember picture,overlay]
% Background image
    \node[above right,inner sep=0pt] at (current page.south west)
    {
        \includegraphics[width=\paperwidth]{img/blueCS_background.png}
    };
% Logo CS
\node[below right =6cm and 0.5cm] at (current page.north west)
{
    \includegraphics[width=2cm]{img/CS_white_logo.png}
};
\draw[yellowCS, ultra thick] (2,-3.25) -- (8,-3.25);
  
% section title
\node[color=yellowCS,below right=2cm and 2cm] {\Huge \sectionTitle};
\node[color=white,below right=3.5cm and 2cm] {\Large \sectionSubtitle};

\end{tikzpicture}
    
\end{frame}


\begin{frame}{Vulnerability Summary}
    \textbf{Defensive Analysis (Snyk):}
    \begin{itemize}
        \item SQL Injection in \texttt{app.py} (High)
        \item Improper Certificate Validation (High)
        \item DOM-based XSS (Medium)
    \end{itemize}
    \vspace{0.5cm}
    \textbf{Offensive Analysis :}
    \begin{itemize}
\item (Burp Suite) \color{red}\textbf{Critical SQL Injection in /api/transactions}
\item API Documentation Exposure (/api/docs)
\item Nessus: PlainText Credentials
\item Clickbaiting.
    \end{itemize}
\end{frame}

\begin{frame}{Exploitation: SQL Injection (Manual)}
    \begin{itemize}
        \item \textbf{Endpoint:} \texttt{GET /api/transactions}
        \item \textbf{Method:} Intercepted request via Burp Suite Proxy.
        \item \textbf{Payload:} \texttt{' OR '1'='1}
        \item \textbf{Result:} The server returned the transaction history for \textbf{ALL} users.
    \end{itemize}
    \begin{figure}[h]
        \centering
        % Asigura-te ca imaginea e in folderul Figures
        \includegraphics[height=0.45\textheight]{Figures/api_burpsuite_repeater.png}
        \caption{Burp Suite Repeater Proof of Concept}
    \end{figure}
\end{frame}


\begin{frame}{Automatic Plaintext exploit}
    \begin{itemize}
        \item Set up MitM
        \item Intercept.
        \item Profit
    \end{itemize}
    \begin{lstlisting}[language=Bash]
    > uv run proxy.py
    [*] Proxy Server Started, listening on 0.0.0.0:8080 (everywhere)
    [*] But local IP Address is 192.168.0.94:8080
    so Configure your second device to connect to that as an HTTP Proxy
    
    
    [+] INTERCEPTED POST REQUEST TO: vulnbank.org
        [!] CREDENTIALS FOUND IN BODY:
        {"username":"suuup dude","password":"sper ca nu vezi aasta lol lmao that woild be so emberssaing"}
    --------------------------------------------------
    \end{lstlisting}
\end{frame}


\begin{frame}{Clickbaiting exploit}
    \begin{itemize}
        \item Load transparent Iframe that contains some logging in info, or anything important that the user doesn't neccesarily want to click.
        \item Load something unassuming below it
        \item Profit
    \begin{figure}[h]
        \centering
        \includegraphics[height=0.45\textheight]{Figures/theghostly-figure.png}
        \caption{Ghostly iframe}
    \end{figure}
    \end{itemize}
\end{frame}


% --- SECTIUNEA 4: REMEDIATION ---
\renewcommand*\sectionTitle{Remediation}
\renewcommand*\sectionSubtitle{Fixing the Issues}
% Section 2 slide frame
\begin{frame}[plain]
\begin{tikzpicture}[remember picture,overlay]
% Background image
    \node[above right,inner sep=0pt] at (current page.south west)
    {
        \includegraphics[width=\paperwidth]{img/blueCS_background.png}
    };
% Logo CS
\node[below right =6cm and 0.5cm] at (current page.north west)
{
    \includegraphics[width=2cm]{img/CS_white_logo.png}
};
\draw[yellowCS, ultra thick] (2,-3.25) -- (8,-3.25);
  
% section title
\node[color=yellowCS,below right=2cm and 2cm] {\Huge \sectionTitle};
\node[color=white,below right=3.5cm and 2cm] {\Large \sectionSubtitle};

\end{tikzpicture}
    
\end{frame}


\begin{frame}{Remediation Steps & Reporting}
    \begin{block}{1. Fixing SQL Injection}
    Replaced dynamic SQL concatenation with \textbf{Parameterized Queries} (Prepared Statements).
    \end{block}
    
    \begin{block}{2. Hardening API}
    Disabled Swagger UI in production environment and enforced Authentication on docs endpoints.
    \end{block}
    
    \begin{block}{3. Reporting (RIMSEC)}
    Reported as \textbf{High Severity}.
    \begin{itemize}
        \item \textbf{R}eproducible: 100\% via Burp.
        \item \textbf{I}mpact: Critical Data Leak.
        \item \textbf{M}itigation: Code patch applied.
    \end{itemize}
    \end{block}
\end{frame}

% --- CONCLUSIONS ---
\begin{frame}{Conclusions}
\begin{itemize}
    \item \textbf{Synergy:} The Defensive approach (SAST) identified potential flaws, while the Offensive approach (DAST) validated their real-world impact.
    \item \textbf{Tooling:} Automated tools (Nessus/Snyk) are excellent for baselines, but manual testing (Burp Suite) is essential for logic flaws like the API SQL Injection.
    \item \textbf{Lesson Learned:} Security must be layered. A simple configuration error (exposed docs) led to a critical data breach.
\end{itemize}
\end{frame}

% --- CONTACT PAGE ---
% Contact slide frame
\begin{frame}[plain]
\begin{tikzpicture}[remember picture,overlay]
% Background image
\node[above right,inner sep=0pt] at (current page.south west)
{
    \includegraphics[width=\paperwidth]{img/blueCS_background.png}
};
    
% Facultate
\node[color=white,below left=-0.7cm and -11cm]{ \small \MakeUppercase\Facultate};
% Universitate
\node[color=white,below left=-0.25cm and -9.5cm]{\small \MakeUppercase\Universitate};

% adresa
\node[color=white,below left=0.5cm and -9.9cm,align=center,text width=4cm]{\footnotesize\Adresa};

% website
\node[color=white,below left=1.75cm and -8.9cm,align=center]{\small \Website};
% Logo CS
\node[below right=3cm and 1cm] at (current page.north west)
{
    \includegraphics[width=3cm]{img/CS_white_logo.png}
};
\draw[white, ultra thick] (4,0.65) -- (4,-2.15);

\end{tikzpicture}
    
\end{frame}


\end{document}